\chapter{Cơ sở lý thuyết \& Các công trình liên quan}
\label{chap:co-so-ly-thuyet}

Chương này sẽ trình bày hai nền tảng lý thuyết cốt lõi của nghiên cứu là Thang phân loại Bloom và Hệ thống đa tác tử. Ngoài ra, báo cáo sẽ giới thiệu về thư viện LangGraph được sử dụng để triển khai kiến trúc đề xuất kèm theo kỹ thuật sử dụng LLM để đánh giá kết quả (LLM-as-a-Judge). Phần cuối sẽ đi sâu vào các phương pháp hiện tại và chỉ ra những hạn chế cũng như các khoảng trống nghiên cứu còn tồn tại.

\section{Phân loại Bloom}

Để các hệ thống sinh câu hỏi trắc nghiệm tự động có thể vừa đảm bảo các mục tiêu sư phạm vừa không biến thành một cỗ máy tạo câu hỏi ngẫu nhiên, việc thiết kế độ khó cần dựa trên một thang đo đã được kiểm chứng bởi khoa học. Nghiên cứu này tập trung vào khung nhận thức được công nhận rộng rãi nhất trong lĩnh vực giáo dục là \textbf{phân loại Bloom}. Thang đo này được nhà tâm lý học và giáo dục \textit{Benjamin Bloom} đề xuất lần đầu vào năm 1956 \cite{bloom1956taxonomy} và được tinh chỉnh bởi \textit{Anderson} và \textit{Krathwohl} vào năm 2001 \cite{anderson2001taxonomy}. Trong nghiên cứu này, nhằm tối ưu hóa và phù hợp với hình thức thi câu hỏi trắc nghiệm đa lựa chọn (MCQ - Multiple-Choice Question), sáu cấp độ ban đầu của thang đo Bloom được phân loại lại thành ba cấp độ cốt lõi là Truy xuất, Suy luận và tổng hợp, và Lập luận phản biện.

\subsection{Cấp độ 1 - Truy xuất}

Cấp độ Truy xuất tương ứng với bậc Nhớ trong thang đo Bloom sửa đổi. Cấp độ này là nền tảng của sự nhận thức, tập trung chủ yếu vào khả năng ghi nhớ cơ học và truy xuất thông tin dài hạn, hoặc nhận dạng một đối tượng, sự kiện được trình bày rõ ràng trong một đoạn văn bản thuộc bài đọc. Quá trình này thường được gọi là nhận thức nông.

Trong quá trình tinh chỉnh các chỉ dẫn cho LLMs sinh câu hỏi ở cấp độ Truy xuất, hệ thống được cài đặt để tập trung vào các chi tiết cụ thể, định nghĩa thuật ngữ và nhận dạng các sự kiện. Đặc thù của MCQ ở dạng này thường xoay quanh các câu hỏi “Ai?”, “Cái gì?”, “Ở đâu?” và “Khi nào?”. Đối với bài toán  này, LLMs thể hiện vô cùng xuất sắc trong việc tự động tạo ra các câu hỏi, bởi vì quy trình này cơ bản chỉ đòi hỏi khả năng đối sánh chuỗi và trích xuất các thực thể từ văn bản mà không cần hiểu sâu hay xây dựng bất kỳ chuỗi suy luận logic nào. Không chỉ vậy, các phương án nhiễu ở mức độ này cũng tương đối dễ xây dựng; chúng thường là các thực thể hoặc dữ liệu xuất hiện trong văn bản gốc để tạo ra sự khó khăn cho người đọc.

\subsection{Cấp độ 2 - Suy luận và tổng hợp}

Cấp bậc thứ hai bao gồm hai kỹ năng Suy luận và Tổng hợp, đại diện cho các bậc Hiểu, Áp dụng và Sáng tạo trong phân loại Bloom. Ở cấp độ này, thay vì được hỏi các câu hỏi có thể truy xuất trực tiếp từ bề mặt của văn bản, người học sẽ buộc phải xử lý thông tin đã học và ứng dụng lý thuyết đó vào một ngữ cảnh hoàn toàn mới hoặc một tình huống giả định.

Trong bối cảnh bài toán MCQ, Suy luận yêu cầu người học phải có khả năng liên kết các mệnh đề nằm rải rác trong bài đọc để rút ra các kết quả, nhận định không được tác giả viết ra một cách trực tiếp. Đối với Tổng hợp, cấp độ này yêu cầu người học phải chắt lọc, sắp xếp lại các yếu tố riêng lẻ thành một góc nhìn bao quát hoặc dự đoán một kết quả tương lai. Bài toán này đòi hỏi cần có các chỉ dẫn cụ thể và hợp lý cho LLMs. Ngoài ra, việc tạo ra các đáp án nhiễu ở cấp độ này cũng là một thách thức kỹ thuật đối với AI: các đáp án sai phải được thiết kế dựa vào sự suy diễn sai lầm về mặt logic hoặc dựa vào việc người đọc hiểu sai một quy luật nào đó để từ đó có thể tạo ra các lựa chọn bẫy hợp lý về mặt trực giác. 

\subsection{Cấp độ 3 - Lập luận phản biện}

Cấp độ cuối cùng của nghiên cứu này là Cấp độ Lập luận phản biện, tương ứng với bậc Phân tích và Đánh giá. Trọng tâm của cấp độ nhận thức này là khả năng chia nhỏ một lượng thông tin đồ sộ, phức tạp thành các thành phần nhỏ hơn nhằm phân tích, nhận diện các giả định ngầm và đưa ra những nhận xét, quan điểm khách quan dựa trên sự đối chiếu với các bằng chứng thực tiễn.

Các câu hỏi trắc nghiệm ở cấp độ này thường yêu cầu người học tưởng tượng ra các tình huống kiểm tra ngụy biện logic, yêu cầu họ phân biệt rõ ràng giữa một sự thật khách quan và một quan điểm chủ quan, hoặc được yêu cầu đánh giá độ chính xác, hợp lý của một lập luận khoa học. Thách thức lớn nhất mà các LLM phải đối mặt ở cấp độ này là việc sinh ra các câu hỏi chất lượng cao chỉ có duy nhất một đáp án đúng. Quan trọng hơn cả, phương án nhiễu ở mức độ này không thể chỉ là những thông tin sai lệch thông thường, mà phải là các lập luận, lý lẽ nhìn bề mặt có vẻ cực kỳ thuyết phục nhưng lại chứa các lỗ hổng về mặt logic, từ đó buộc người học phải sử dụng tư duy phản biện cao mới có thể nhận diện được.

\section{Hệ thống đa tác tử}

\subsection{Khái niệm hệ thống đa tác tử}

Kiến trúc đa tác tử là một kỹ thuật hiện đại cho phép các LLM có thể xử lý các bài toán phức tạp mà một LLM duy khó có thể xử lý một cách chính xác và hiệu quả. Nguyên lý cốt lõi của kỹ thuật này là chia nhỏ bài toán ban đầu thành nhiều nhiệm vụ nhỏ hơn và giao cho một mạng lưới các tác tử xử lý. Trong hệ thống này, mỗi tác tử sẽ được chỉ định thực hiện một nhiệm vụ duy nhất. Các tác tử có thể tương tác qua lại với nhau để trao đổi thông tin. Triết lý này khác với việc yêu cầu một LLM giải quyết trọn vẹn một bài toán, một công việc mà đòi hỏi quy trình phức tạp về cả logic và nghiệp vụ, khiến kết quả đầu ra kém chất lượng và có nguy cơ ``ảo giác’’ (hallucination) cao.

Mỗi tác tử trong kiến trúc này được thiết kế để mô phỏng cách mà con người phân công và xử lý công việc. Kỹ thuật \textbf{ReAct} \cite{yao2022react} thường được sử dụng để yêu cầu các LLM thực hiện một chuỗi hành động lặp lại gồm suy nghĩ, hành động, quan sát kết quả và điều trình suy nghĩ mỗi khi thực hiện một tác vụ. Nhằm tối ưu tài nguyên nhân thức của các LLM, đặc biệt là các sLLM, mỗi tác tử thường sẽ nhận được một chỉ dẫn hẹp với nhiệm vụ hoàn thành một công việc duy nhất. Ví dụ, Tác tử \italic{Sinh câu hỏi} chỉ có nhiệm vụ tạo ra các câu hỏi trắc nghiệm với một cấp độ Bloom nhất định trong khi tác tử \italic{Phân loại} lại có nhiệm đánh giá và dự đoán cấp độ của các câu hỏi vừa tạo, từ đó giúp tác tử {Sinh câu hỏi} có thể điều chỉnh cho phù hợp. Sự kết hợp giữa tính chuyên môn hóa cao và sự tương tác giữa các tác tử giúp cải thiện chất lượng nội dung sinh ra, duy trì độ ổn định và giảm hiện tượng ảo giác thường thấy ở các LLM.

\subsection{Thư viện LangGraph}

Nghiên cứu sử dụng thư viện LangGraph để triển khai và cài đặt kiến trúc đa tác từ đề xuất  một cách đúng đắn và thuận tiện nhất. LangGraph là một thư viện mã nguyên mở chuyên dụng dành cho phát triển các hệ thống LLM và là một thành phần quan trọng trong hệ sinh thái LangChain nổi tiếng. Tuy nhiên, khác với triết lý tuyến tính của LangChain truyền thống, LangGraph sử dụng một đồ thị có hướng để điều phối các tác tử kèm theo một đối tượng trạng thái để lưu trữ dữ liệu. Ba thành phần cốt lõi của một hệ thống đa tác tử khi được triển khai bằng thư viện LangGraph là:

Trạng thái: Đây là một đối tượng lưu trữ dữ liệu tập trung của toàn bộ đồ thị. Nó hoạt động như một bộ nhớ chung nhằm duy trì ngữ cảnh cho hệ thống thống. Một số loại dữ liệu cần lưu trong cấu trúc dữ liệu này có thể kể đến như danh sách các câu hỏi đã tạo và lỗi của từng câu.

Nút: Đây là nơi dữ liệu được xử lý thông qua các hàm Python. Logic và nghiệp vụ của hệ thống sẽ được triển khai trong các đối tượng Nút này. Quy trình hoạt động bắt đầu từ bước tiếp nhận trạng thái hiện tại, gửi lệnh đến LLM để xử lý và kết thúc bằng cách cập nhật đầu ra vào đối tượng Trạng thái chung.

Cạnh: Thành phần này có nhiệm vụ xác định hướng di chuyển của luồng dữ liệu giữa các nút của đồ thị có hương. Mỗi cạnh có thể bao gồm các điều kiện rẽ nhánh đặc biệt hoặc có khả năng kích hoạt cơ chế thực thi song song. Cạnh sẽ được nối với Nút `END` mô tả trạng thái kết thúc của luồng dữ liệu.

Nhờ cách thiết kế trên, thư viện LangGraph cho phép các hệ thống đa tác tử có khả năng tự định hướng luồng dữ liệu. Nếu một tác tử \italic{Phân loại} phát hiện ra một câu hỏi sinh ra có cấp độ nhận thức không đúng so với cấu hình, tác tử này có thể yêu cầu tác tử {Sinh câu hỏi} tạo lại câu hỏi đó. Tương tự, nếu tác tử {Học sinh} phát hiện ra một câu hỏi có các đáp án bị lỗi, nó hoàn toàn có thể yêu cầu sửa đổi câu hỏi thông qua tác tử {Hiệu chỉnh}. Các vòng lặp này chỉ dừng lại khi các tác tử kiểm tra chấp nhận kết quả sinh.

\section{Phương pháp đánh giá bằng LLM}

Việc đánh giá chất lượng của dữ liệu tạo sinh luôn là một bài toán phức tạp trong lĩnh vực trí tuệ nhân tạo, đặc biệt là xử lý ngôn ngữ tự nhiên. Các độ đo truyền thống như BLEU hay ROUGH tỏ ra kém hiệu quả trong việc đánh giá các chất lượng của LLM vì điểm yếu cố hữu của các độ đo này là không hiểu ngữ nghĩa và ngữ cảnh của dữ liệu. Bên cạnh đó, việc đánh giá bởi con người thường tiêu tốn chi phí, thời gian và hạn chế về khả năng mở rộng mặc dù có độ chính xác cao. Phương pháp sử dụng mô hình ngôn ngữ lớn làm giám khảo (LLM-as-a-Judge) giải quyết vấn đề này bằng cách sử dụng chính các LLM thương mại mạnh mẽ để thay thế vai trò của con người trong việc kiểm tra và đánh giá chất lượng dữ liệu đầu ra.

 Liu và các cộng sự đã đề xuất kỹ thuật G-Eval vào năm 2023 để hiện thực hóa  triết lý này. Nguyên lý của G-Eval là sử dụng LLM giám khảo nhận đầu vào gồm 3 thành phần gồm: nội dung nguồn, nội dung cần đánh giá và một bộ tiêu chuẩn đánh giá chi tiết. Đầu ra của LLM giảm khảo có thể là một nhãn nhị phân đúng sai hoặc một điểm số cụ thể. Phương pháp LLM-as-a-Judge mang lại ba lợi thế quan trọng. Thứ nhất là khả năng mở rộng vượt trội. Hệ thống đánh giá có thể xử lý hàng nghìn dữ liệu trong thời gian ngắn và chi phí thấp hơn rất nhiều so với con người. Thứ hai, phương pháp này cung cấp khả năng thay đổi linh hoạt. Bộ tiêu chuẩn có thể được cập nhật và điều chỉnh dễ dàng. Thứ ba, LLM giám khảo có tính ổn định cao hơn con người, kết quả có thể dễ dàng tái lập, loại bỏ được các yếu tố chủ quan của những người đánh giá khác nhau.

Ngoài ra, kỹ thuật này thường được kết hợp với kỹ thuật chuỗi suy luận (CoT - Chain-of-Thought) \cite{} được đề xuất bởi xxxx và các cộng sự vào năm xxxx. Kỹ thuật này yêu cầu LLM giám khảo phải đưa ra các lập luận, đánh giá cụ thể trước khi đưa ra điểm số cuối cùng, nhờ đó giúp cải thiện chất lượng và độ ổn định của quy trình đánh giá.

Tuy nhiên, phương pháp LLM-as-a-Judge vẫn tồn tại một số hạn chế cần được lưu ý. Đầu tiên, LLM giám khảo có thế mắc các lỗi ``thiên vị’’ đối với các kết quả được sinh ra bởi chính LLM cùng loại. Không chỉ vậy, nếu không kết hợp với kỹ thuật như CoT, LLM giám khảo có thể đưa ra các kết luận không chính xác. Chính vì vậy, kết quả từ LLM-as-a-Judge nên được đối chiếu với đánh giá từ con người nhằm xác nhận độ tin cậy.

\section{Các công trình liên quan}

Trong phần này, khóa luận sẽ trình bày các hướng nghiên cứu hiện nay về bài toán sinh câu hỏi trắc nghiệm từ văn bản. Về tổng quan, hai phương pháp sinh câu hỏi chính bao gồm: Phương pháp sinh câu hỏi truyền thống và sinh câu hỏi dựa trên mô hình ngôn ngữ lớn.

\subsection{Sinh câu hỏi trắc nghiệm tự động bằng quy tắc}

Trước khi có sự xuất hiện của các mô hình học sâu, các bài toán liên quan đến xử lý ngôn ngữ tự nhiên (NLP - Natural Language Processing)  thường được giải quyết bằng các quy tắc ngôn ngữ học được thiết kế một cách thủ công bởi chuyên gia. . \textit{Heilman} và \textit{Smith} \cite{heilman2010good} đã đề xuất phương pháp tiếp cận bằng cách xếp hạng thống kế các câu hỏi được sinh ra. Cụ thể, hệ thống sẽ sinh ra một số lượng lớn các câu hỏi từ văn bản đầu vào thông qua các phép biến đổi về mặt cú pháp. Tiếp theo, hệ thống sẽ sử dụng một mô hình xếp hạng thống kê để lọc và chọn ra các câu hỏi có chất lượng sư phạm tốt nhất.

Với một hướng tiếp cận khác, Mitkov và các cộng sự \cite{} lại đề xuất một hệ thống tạo câu hỏi trắc nghiệm sử dụng các kỹ thuật NLP như nhận dạng thực thể có tên, gán nhãn loại từ và phân tích cú pháp. Nhờ đó, hệ thống có thể nhận diện ra các khái niệm, đối tượng quan trọng trong văn bản đầu vào, từ đó có thể áp dụng trực tiếp các quy tắc biển đổi ngôn ngữ để chuyển thành các câu hỏi và các đáp án nhiễu tương ứng.

Mặc dù các phương pháp truyền thống cung cấp cho người dùng khả năng kiểm soát cao về mặt ngữ pháp và cấu trúc câu hỏi, phương pháp này lại có nhiều hạn chế lớn khiến chúng khó có thể ứng dụng trong bài toán thực tế. Thứ nhất, các câu hỏi được tạo ra bằng quy tắc thường có mức độ nhận thức thấp tho thang đo Bloom \cite{}. Các câu hỏi này chủ yếu biến đổi trực tiếp nội dung trong văn bản, điều này dẫn đến việc chỉ có thể sinh ra các câu hỏi tái hiện thông tin bề mặt, chất lượng sư phạm không cao. Thứ hai, các phương án nhiễu của phương pháp này không gây khó khăn cho học sinh, thường bị loại trừ một cách dễ dàng. Thứ ba, công việc thiết kế ra các bộ quy tắc ngôn ngữ cho hệ thống sinh đòi hỏi chuyên môn cao về ngôn ngữ học và tiêu tốn thời gian và tài nguyên.

\subsection{Sinh câu hỏi trắc nghiệm dựa trên mô hình ngôn ngữ lớn}

Sự ra đời của các LLM đã mở ra một hướng đi hoàn toàn mới trong lĩnh vực tạo sinh của ngành giáo dục, đặc biệt là với bài toán tạo câu hỏi trắc nghiệm từ văn bản. Với khả năng tạo ra các văn bản của cấu trúc ngữ pháp và từ vựng chính xác, các hệ thống sinh có thể loại bỏ hoàn toàn các quy tắc ngữ pháp thủ công phức tạp. Không những thế, các LLM còn có khả năng hiểu ngữ cảnh và chỉ dẫn một cách chặt chẽ.

\textit{Elkins} và công sự \cite đã thực hiện các đánh giá chất lượng của các câu hỏi được sinh ra bởi các LLM trong lĩnh vực giáo dục. Kết quả nghiên cứu cho thấy mặc dù các mô hình ngôn ngữ này có thể tạo ra các câu hỏi có chất lượng đạt yêu cầu để đưa vào ứng dụng trong giảng dạy thực tế, mức độ nhận thức của các câu hỏi này chưa cao theo thang đo Bloom. Nghiên cứu này chỉ ra rằng khoảng cách giữa các câu hỏi ghi nhớ bề mặt đơn thuần và câu hỏi yêu cầu phân tích, suy luận vẫn là một thách thức lớn đối với các LLM hiện nay.

\section{Phân tích khoảng trống nghiên cứu}

Từ quá trình tìm hiểu và tổng quát hóa các công nghiên nghiên cứu liên quan, khóa luận rút ra hai vấn đề hiện hữu của bài toán sinh câu hỏi trắc nghiệm tự động:

\bold{Khả năng kiểm soát sư phạm và độ chính xác thấp ở tư duy bậc cao.} Các phương pháp sinh câu hỏi bằng quy tắc truyền thống như phân tích cú pháp bằng NLP hay từ điển WordNet tỏ ra cứng nhắc, không linh hoạt và yêu cầu công sức con người để thiết kế. Mặt khác, việc sử dụng trực tiếp một LLM để tạo ra toàn bộ một bài kiểm tra trắc nghiệm mà thiếu đi các bước đánh giá thường dẫn đến kết quả có chất lượng thấp hoặc gặp tình trạng ảo giác (halluciation).

\bold{Tốn kém tài nguyên và rào cản triển khai thực tế.} Các hệ thống giáo dục hiện nay thường phụ thuộc quá nhiều vào các LLM thương mại. Mặc dù các mô hình này đem lại chất lượng câu hỏi tương đối tốt, chúng thường đi kèm với chi phí đắt đỏ và độ trễ cao. Ngoài ra, vấn đề bảo mật dữ liệu cũng cần được xem xét, đặc biệt là đối với các hệ thống học liệu bản quyền của các tổ chức giáo dục lớn.





